\documentclass[10pt,a4paper]{amsart}
\usepackage[margin=19mm]{geometry}
\usepackage{amsmath,amssymb,mathtools}
\usepackage[scr=rsfs,cal=euler]{mathalfa}
\usepackage{xfrac}
\usepackage[unicode,hidelinks,bookmarks=false]{hyperref}
\newcommand{\ie}{i.e.\ }
\newcommand{\cf}{cf.\ }
\newcommand{\resp}{respectively\ }
\newcommand{\Subsection}{\subsection}
\providecommand{\nobreakdash}{\nobreak\hbox{-}}
\setlength{\emergencystretch}{2em}
\multlinegap0pt
\usepackage{mathtools}
\usepackage{xcolor}
\usepackage{amssymb}
\usepackage{bm}
\newtheorem{lem}{Lemma}
\def \R {{\mathbb R}}
\newcommand{\SL}{\operatorname{SL}}
\newcommand{\SO}{\operatorname{SO}}
\def \Z {{\mathbb Z}}
\newcommand{\vol}{\operatorname{vol}}
\newcommand{\x}{\mathbf{x}}
\newcommand{\y}{\mathbf{y}}
\title{The divisor function for matrices}
\author{Tim Browning, Nikita P. Kalinin, Alina Ostafe, Damaris Schindler, Lena Wurzinger}
\date{Source: arXiv:2608.25078v1 (CC BY 4.0)}
\begin{document}
\maketitle

\noindent\emph{Source and licence.} The divisor function for matrices. Version arXiv:2608.25078v1 (CC BY 4.0), distributed by arXiv under the Creative Commons Attribution 4.0 International licence, http://creativecommons.org/licenses/by/4.0/, as recorded in the arXiv licence metadata for this exact version and shown on the abstract page https://arxiv.org/abs/2608.25078v1. Full licence notice: \texttt{licenses/arxiv-2608.25078-CC-BY-4.0.txt}.\par\smallskip

\noindent\textit{Editorial scope.} Counted unit: the article's lem lem:volumecomputation (source0.tex lines 842--845). The article's complete original proof of it is delivered with it (source0.tex lines 846--1163, one \texttt{proof} environment of 318 lines). No mathematics is written, repaired or re-derived: the statement and the proof are the article's own lines, in the article's own notation, and no retained line is substituted or re-flowed.  Retained uncounted as the source-local context the proof needs: lines 578--584; lines 633--636; lines 637--645; lines 649--661; lines 673--676.  Nothing was omitted from any retained span, so the package carries no omission record.  No retained line contains a citation, so the wrapper carries no bibliography and no bibliography entry.  No retained line contains a figure environment, an image-inclusion command or drawing code, so no image is delivered and the package declares no asset dependency.  Editorial formatting only, in addition: an AMS \texttt{amsart} preamble that restates the article's own declarations and macros used by the retained lines, namely the environments \texttt{lem} on separate counters, together with the standard packages those lines need.  The retained environments print in this document as Lemma 1 in order of appearance, whereas the article prints the same items with the numbering of its own full document.  The complete native source of the article as distributed in the arXiv e-print for this exact version is bundled unchanged as \texttt{sources/208542/source0.tex}.
\begin{align}\label{eq:basic_infty_norm_inequ}
    | M + L | \le  | M| +
    | L| \quad
    \text{ and } \quad
    | M  L| \le  n| M | 
    | L|.
\end{align}

    \begin{align}
    \label{eq:cartan_decomp}
         k_1, k_2 \in \SO_n(\R) =:K, \quad a_{\mathbf{t}} = \mathrm{diag}(e^{t_1}, \ldots, e^{t_n}), \quad \mathbf{t} = (t_1,\ldots, t_n) \in H^{+}, 
    \end{align}

    with     \begin{align}\label{eq:defH+}
  H^{+} := 
 \left \{
  \mathbf{t} \in \R^n : t_1 \ge \cdots \ge t_n , \ \sum_{i=1}^n t_i = 0
 \right  \} \ \subseteq \
 \left \{ 
         \mathbf{t} \in \R^n : \sum_{i=1}^n t_i = 0
         \right \} =: H .
     \end{align} 

  \begin{align}\label{eq:cartanvol}
  c_n \prod_{i<j} \sinh(t_i - t_j) \;
      \mathrm{d}k_1
      \mathrm{d} \mathbf{t}
      \mathrm{d}k_2
      , \qquad
      c_n =  
      \frac{
      2^{\binom{n}{2}+n-1} \pi^{\frac{n(n+1)}{2}-1}
      }{ \prod_{j=2}^n
       (\Gamma(\tfrac{j}{2})^{2} \zeta(j))
      },
  \end{align} 

\begin{align}\label{def:kappa}
    \kappa(X,Y) = \lim_{T \to \infty} 
   T^{-\lfloor n^2/2 \rfloor} \vol (E_{X,Y}(T)),
\end{align}

\begin{lem}\label{lem:volumecomputation}
    For fixed invertible matrices $X,Y \in \mathcal{M}_{n}(\Z)$, the limit $\kappa (X,Y)$ in
    \eqref{def:kappa} exists and satisfies $\kappa (X,Y)>0$.
\end{lem}

\begin{proof}
We use the 
Cartan decomposition     $g = k_1 a_{\mathbf{t}} k_2$ of any 
 $g \in \SL_n(\R)$, in the notation of \eqref{eq:cartan_decomp}.
   Set $T = e^L$ and write
  \begin{align*} 
  f_L(k_1,\mathbf{t},k_2) &= \mathbf{1}_{[| X k_1 a_{\mathbf{t}} k_2 | \le T]} 
      \mathbf{1}_{[| (k_1 a_{\mathbf{t}} k_2)^{-1} Y | \le T]}  \\
      &=
      \mathbf{1}_{[| e^{-L} X k_1 a_{\mathbf{t}} k_2 | \le 1]}  
      \mathbf{1}_{[|e^{-L} (k_1 a_{\mathbf{t}} k_2)^{-1} Y | \le 1]} .
  \end{align*}
  By \eqref{eq:cartanvol}, we have  
  \begin{align*}
      \vol(E_{X,Y}(T)) = c_n \int_K \int_{H^{+}}  \int_K f_L(k_1,\mathbf{t},k_2) \prod_{i<j} \sinh(t_i - t_j) \;
      \mathrm{d}k_1
      \mathrm{d} \mathbf{t}
      \mathrm{d}k_2.
  \end{align*}
  To this integral we apply the change of variables $\mathbf{x} = \mathbf{t} - L\mathbf{v}$, where we set 
  \begin{align}\label{eq:definitionof v}
  \mathbf{v} =
\begin{cases}
(\underbrace{1,\dots,1}_{n/2},
 \underbrace{-1,\dots,-1}_{n/2}) &\text{if $n$ is even},\\
(\underbrace{1,\dots,1}_{(n-1)/2},
 0,
 \underbrace{-1,\dots,-1}_{(n-1)/2}) & \text{if $n$ is odd.}
\end{cases}
  \end{align}
  We further set  $\mathcal{D}_L = H^+ -L \mathbf{v}$. It follows that the volume of $E_{X,Y}(T)$ may be computed as
  \begin{align}\label{eq:volEafterchangeofvariables}
      c_n \int_K \int_{H}  \int_K 
     \mathbf{1}_{\mathcal{D}_L}(\mathbf{x})
     f_L(k_1,\mathbf{x}+L\mathbf{v},k_2) \prod_{i<j} \sinh(x_i - x_j + L(v_i-v_j)) \;
      \mathrm{d}k_1
      \mathrm{d} \x
      \mathrm{d}k_2,
  \end{align}
  where $H$ is given in \eqref{eq:defH+}. 
  We now study the pointwise behaviour of the integrand as $L \to \infty$. Ultimately, via an application of dominated convergence, this will allow us to describe the asymptotic behaviour of the integral.

\subsubsection*{Asymptotics of $\mathbf{1}_{\mathcal{D}_L}$}
    Let us fix any $\mathbf{x} \in H$. We have $\mathbf{x} \in \mathcal{D}_L = H^{+}-L\mathbf{v}$ if and only if 
    \[
        x_i + L v_i
        \ge 
        x_j + L v_j
    \]
    for all $i < j$; therefore, we have
    \begin{align}\label{eq:domoconDL}
        \mathcal{D}_L  \subset \mathcal{C}  :=
    \{
    \mathbf{y} \in H: y_i \ge y_j \text{ if $i < j \le \tfrac{n}{2}$ or $\tfrac{n+1}{2} < i <j$ } \}.
    \end{align}
     Conversely, for $\y \in \mathcal{C}$, we have $y_i + Lv_i \ge y_j + L v_j$ for all $i<j$;
      i.e.\ $\y \in \mathcal{D}_{L}$, 
     for all sufficiently large $L$. Since we study the pointwise behaviour of $\mathbf{1}_{\mathcal{D}_{L}}$, the lower bound on $L$ is allowed to depend on $\y$. 
     Indeed, for $i<j$ with $i < j \le \tfrac{n}{2}$ or $\tfrac{n+1}{2} < i <j$, we have $v_i = v_j$ and $y_i \ge y_j$, and so
      $y_i + L v_i \ge y_j + Lv_j$ holds for all $L$. For the remaining combinations of indices $i<j$, we have  $v_i > v_j$ and so we get $y_i + L v_i \ge y_j + Lv_j$  for all sufficiently large $L$.
    For all $\mathbf{x} \in H$, 
     the pointwise convergence of the indicator function
    \[
    \mathbf{1}_{\mathcal{D}_L}(\x) \to \mathbf{1}_{\mathcal{C}}(\x), 
    \]
    follows as $L \to \infty$.

\subsubsection*{Asymptotics of $f_L$} Recall the notation~\eqref{eq:cartan_decomp}. 
   We observe, as $L \to \infty$, that 
      \begin{equation}\label{eq:defofD+}
          e^{-L} a_{\mathbf{x}+L\mathbf{v}} = \mathrm{diag}(
      e^{x_i + (v_i - 1)L }
      )_{i=1,\ldots,n} \to \mathrm{diag}
      (
      \mathbf{1}_{[v_i=1]} e^{x_i  }
      )_{i=1,\ldots,n} =: D_+(\mathbf{x}) 
      \end{equation}
      and
        \begin{equation}\label{eq:defofD-}
      e^{-L} a_{\mathbf{x}+L\mathbf{v}}^{-1} = \mathrm{diag}(
      e^{-x_i - (v_i + 1)L }
      )_{i=1,\ldots,n} \to \mathrm{diag}
      (
      \mathbf{1}_{[v_i=-1]} e^{-x_i  }
      )_{i=1,\ldots,n} =: D_{-}(\mathbf{x}).
        \end{equation}
      Hence, for all $(k_1,\mathbf{x},k_2)$ not lying on $| X k_1 D_+(\mathbf{x}) k_2|=1$ or $|  k_2^{-1} D_-(\mathbf{x}) k_1^{-1} Y |=1$, we have  
     \begin{equation}
     \label{eq:f(k_1,x,k_2)}
       f_L(k_1,\mathbf{x}+L\mathbf{v},k_2) 
       \to \mathbf{1}_{[| X k_1 D_+(\mathbf{x}) k_2|\le 1]} 
        \mathbf{1}_{[|  k_2^{-1} D_-(\mathbf{x}) k_1^{-1} Y |\le 1]} =: f(k_1,\mathbf{x},k_2) ,
    \end{equation}
      as $L \to \infty$.
      
    For $(k_1,\mathbf{x},k_2)$ lying on $| X k_1 D_+(\mathbf{x}) k_2|=1$ or $|  k_2^{-1} D_-(\mathbf{x}) k_1^{-1} Y |=1$, pointwise convergence 
    of $f_L$ to $f$ 
    may fail: for instance, we may have
    that $\lim_{L \to \infty} | e^{-L} X k_1 a_{\mathbf{x}+L\mathbf{v}} k_2 | =1$ but $| e^{-L} X k_1 a_{\mathbf{x}+L\mathbf{v}} k_2 | > 1 $ for all $L$.
However, the set         \[
       \{(k_1,\x,k_2) :
        | X k_1 D_+(\mathbf{x}) k_2| = 1
        \text{ or }
        |  k_2^{-1} D_-(\mathbf{x}) k_1^{-1} Y | =1
        \}
        \]
        has measure zero.
       To formally justify this, one may cover $K \times K$ by finitely many real 
       analytic coordinate charts. On
each such chart, each entry of the matrices
$Xk_1D_+(\mathbf{x})k_2$ and $k_2^{-1}D_-(\mathbf{x})k_1^{-1}Y$ is real
analytic in $(k_1,\mathbf x,k_2)$. None of these entries is identically equal
to $\pm 1$, as the entries of $D_+(\mathbf x)$ and
$D_-(\mathbf x)$, and thus the entries of $Xk_1D_+(\mathbf{x})k_2$ and $k_2^{-1}D_-(\mathbf{x})k_1^{-1}Y$, can be made arbitrarily small. The vanishing loci of real-analytic functions that are not identically zero, such as $(X k_1 D_+(\mathbf{x}) k_2)_{i,j}\pm1$ and $(k_2^{-1} D_-(\mathbf{x}) k_1^{-1} Y)_{i,j}\pm1$, always constitute a set of measure zero, and our exceptional set is contained in the finite union over $1 \le i,j \le n$ of these sets.
Hence we have pointwise almost everywhere convergence 
$f_L \to f$.

      
      Moreover, by \eqref{eq:basic_infty_norm_inequ} and the fact that  $| k | \le 1$ for all $k \in K$,  for all $\x \in \mathcal{D}_{L}=H^{+}-L\mathbf{v}$ with $| e^{-L} X k_1 a_{\mathbf{t}} k_2 | \le 1$  and $|e^{-L} (k_1 a_{\mathbf{t}} k_2)^{-1} Y | \le 1$, writing $\mathbf{t}=\mathbf{x}+L\mathbf{v}$, we have 
      \begin{align*}
           e^{x_1} = 
      e^{-L}
      | a_{\mathbf{x} + L\mathbf{v}} | 
      &=
      e^{-L}| (X k_1)^{-1}X k_1 a_{\mathbf{x} + L\mathbf{v}} k_2 k_2^{-1} | \\
      &\le
      n^3 |  X^{-1} |  | e^{-L} X k_1 a_{\mathbf{x} + L\mathbf{v}} k_2  | \le n^3 |  X^{-1} |
      \end{align*}
      and 
      \[
      e^{-x_n} = 
      e^{-L} |a_{\mathbf{-x} - L\mathbf{v}} | 
      % = 
      % | k_2 (k_1 a_{\mathbf{t}} k_2)^{-1} Y Y^{-1} k_1| 
      \le n^3
      | Y^{-1} |
      | e^{-L}(k_1 a_{\mathbf{t}} k_2)^{-1} Y  | \le  n^3
      | Y^{-1} |.
      \]
       By continuity, analogous bounds hold for $|D_+(\x)|$ and 
      $|D_-(\x)|$. Namely, if we have  $\limsup_{L \to \infty}| e^{-L} X k_1 a_{\mathbf{t}} k_2 | \le 1$  and $\limsup_{L \to \infty}|e^{-L} (k_1 a_{\mathbf{t}} k_2)^{-1} Y | \le 1$, we get 
      \[
      e^{x_1}=
      |D_+(\x)| = \lim_{L\to \infty} ( e^{-L} |a_{\mathbf{x}+  L\mathbf{v}}|) \le n^3|X^{-1}|
      \]
      and 
      \[
      e^{-x_n}=
       |D_-(\x)| = \lim_{L\to \infty} ( e^{-L} |a^{-1}_{\mathbf{x}+  L\mathbf{v}}|) \le n^3|Y^{-1}|.
      \]
    Taking logarithms, it follows that there exists a constant $c_0 = c_0(X,Y) > 0$, independent of $L$, such that 
    \begin{align} \label{eq:suppboundforf}
        x_1 \le c_0, \ x_n \ge - c_0 
    \end{align}
    whenever $\mathbf{1}_{\mathcal{D}_{L}}(\x)f_L(k_1,\mathbf x+L\mathbf{v},k_2)\neq0$ or
$\mathbf{1}_{\mathcal{C}}(\x)f(k_1,\mathbf x,k_2)\neq0$, where $f(k_1,\mathbf x,k_2)$ is defined by~\eqref{eq:f(k_1,x,k_2)}.


\subsubsection*{Asymptotics of $\prod_{i<j} \sinh(x_i-x_j + L(v_i-v_j))$} 
For $v_i > v_j $, since $\sinh(x) = \frac{1}{2}(e^x - e^{-x})$, we see that 
      \begin{align}\label{eq:sinhasymp}
          \sinh(x_i - x_j + L(v_i-v_j)) 
       \sim \frac{e^{x_i - x_j + L(v_i-v_j)}}{2}, \quad \text{as $L \to \infty$,}
      \end{align}
            while for $v_i = v_j$, we have
      \[
      \sinh(x_i - x_j + L(v_i-v_j)) = \sinh(x_i-x_j).
      \]
      Since $\sum_{i<j}(v_i-v_j) = \lfloor n^2/2 \rfloor$, this yields, as $L \to \infty$,   
      \[
      \prod_{i<j} \sinh(x_i - x_j + L(v_i-v_j)) \sim 
      e^{\lfloor n^2/2 \rfloor L} \Psi(\mathbf{x})
      ,
      \]
      where 
      \[
      \Psi(\mathbf{x}) = 2^{-\binom{n}{2}+2\binom{\lfloor n/2 \rfloor}{2}} \prod_{\substack{i < j \\ v_i = v_j}}
      \sinh(x_i-x_j)
      \prod_{\substack{i < j \\ v_i > v_j}} e^{x_i - x_j}.
      \]
    Here we used the fact that there are $\binom{n}{2}-2\binom{\lfloor n/2 \rfloor}{2}$
    factors of the form \eqref{eq:sinhasymp}. Indeed, there
    are $\binom{n}{2}$ indices $i<j$ overall, and of these, there are $2\binom{\lfloor n/2 \rfloor}{2}$ indices  with $v_i = v_j$. As $v_i \ge v_j$ for all $i<j$, it follows that there are $\binom{n}{2}-2\binom{\lfloor n/2 \rfloor}{2}$ indices $i<j$ with $v_i>v_j$.

      
      We observe that for all $\x \in \mathcal{D}_L$, we have
      \begin{align}\label{eq:domconvpsi}
      0 \le
          e^{-\lfloor n^2/2 \rfloor L}\prod_{i<j} \sinh(x_i - x_j + L(v_i-v_j)) \le \prod_{i<j} e^{x_i-x_j} =\exp(F(\x)),
      \end{align}
      where we set $F(\x)=\sum_{i<j}(x_i-x_j)$.

\subsubsection*{Dominated convergence} 
      Let 
      \[
      \mathcal{C}_0 = 
      \{ 
      \y \in \mathcal{C}:  y_1 \le c_0 , y_n  \ge -c_0
      \},
      \]
      with $\mathcal{C}$ defined as in \eqref{eq:domoconDL} and $c_0$ given as in \eqref{eq:suppboundforf}. By \eqref{eq:domoconDL},
      \eqref{eq:suppboundforf} and 
      \eqref{eq:domconvpsi},  the function 
      $\mathbf{1}_{\mathcal{C}_0} (\mathbf{x}) \exp(F(\x))$ gives an upper bound for the modulus of the renormalised integrand from  \eqref{eq:volEafterchangeofvariables}; i.e. for
      \[
       e^{-\lfloor n^2/2\rfloor L}
       \mathbf{1}_{\mathcal{D}_L}(\mathbf{x})
     f_L(k_1,\mathbf{x}+L\mathbf{v},k_2) 
      \prod_{i<j} \sinh(x_i - x_j + L(v_i-v_j)),
      \]
      as well as for the modulus of its almost everywhere pointwise limit 
      $
      \mathbf{1}_{\mathcal{C}}(\x)
     f(k_1,\mathbf{x},k_2)
      \Psi(\x)$.
      
    We now show that there exists a positive constant $\eta > 0$ with
    \begin{align}\label{eq:fboundonR}
         F(\mathbf{r}) \le - \eta | \mathbf{r} |
    \end{align}
    for all $\mathbf{r} \in \mathcal{R}$, where we let 
    \[ \mathcal{R}:= \{\y \in \mathcal{C}: y_1 \le 0, \; y_n \ge 0 \} .
    \]
     From the definition of $\mathcal{C}$ in \eqref{eq:domoconDL}, for $\y \in \mathcal{R}$ we have that
     $0 \ge y_1 \ge y_i$ for all $i$ with $v_i =1$, and  $0 \le y_n \le y_i$ for all $i$ with $v_i =-1$.
    For $\y \in \mathcal{C}_0$, we have $y_1 \le c_0$ and $y_n \ge -c_0$; by the definition of $\mathcal{C}$ in \eqref{eq:domoconDL}, it follows that 
    $c_0 \ge y_1 \ge y_i$ for all $i$ with $v_i =1$, and  $-c_0 \le y_n \le y_i$ for all $i$ with $v_i =-1$. Therefore every $\x \in \mathcal{C}_0$ may be written as 
    \[
    \x = \x_0 + \mathbf{r}, 
    \]
where $|\x_0| \le c_0$ and $\mathbf{r} \in \mathcal{R}$.
    It therefore suffices
    to establish \eqref{eq:fboundonR}  in order to prove that the function $\mathbf{1}_{\mathcal{C}_0} \exp(F(\x))$ is integrable.
    
    Observe that on the polytope 
    \[
    \mathcal{P} = \{ 
    \mathbf{x} \in H^+: x_1 \le 1, x_n \ge -1 
    \},
    \]
    where $H^{+}$ is defined in \eqref{eq:defH+},
    the vector $\mathbf{v}$  in \eqref{eq:definitionof v} is the unique maximiser of the function $F(\x)$. 
    Indeed, writing $d_i = x_{i} - x_{i+1} \ge 0$, we see that 
    \[
    F(\x) = \sum_{i=1}^{n-1} i(n-i)d_i, \qquad \sum_{i=1}^{n-1} d_i = x_1 - x_n \le 2.
    \]
    Thus, considering the coefficients $i(n-i)$, we see that $F(\x)$ is maximised by  choosing $\x$ in such a way that $d_{n/2} = 2$ for $n$ even and $d_{(n-1)/2}+d_{(n+1)/2}=2$ for $n$ odd. Because of the constraint $ x_1 + \cdots + x_n = 0$, this is achieved on $\mathcal{P}$ exactly by $\x=\mathbf{v}$. 

    For every $\mathbf{r} \in \mathcal{R} \setminus \{\mathbf{0}\}$, there is some small positive number $s_\mathbf{r}$ with $\mathbf{v} + s_\mathbf{r}\mathbf{r} \in \mathcal{P}$. By linearity of $F$ and the fact that $\mathbf{v}$ is the unique maximiser, we have that 
    \[
    F(\mathbf{v} + s_\mathbf{r}  \mathbf{r}) = F(\mathbf{v}) + s_\mathbf{r} F(\mathbf{r}) < F(\mathbf{v}), 
    \]
    and so $F(\mathbf{r}) < 0$. As the set $\{ \mathbf{r} \in \mathcal{R}: | \mathbf{r}| =1 \}$ is compact and as $F$ is linear, the desired inequality  \eqref{eq:fboundonR} follows.

      Now it follows by dominated convergence  that $\lim_{T \to \infty}  (T^{-\lfloor n^2 / 2 \rfloor} \vol(E_{X,Y}(T)))$ exists, is finite and may be expressed as
      \begin{align*}
          \lim_{T \to \infty}  (T^{-\lfloor n^2 / 2 \rfloor} \vol(E_{X,Y}(T)))  
          =
          c_n
          \int_{K} \int_{\mathcal{C}} \int_{K}
          G(k_1,\x,k_2)
          \mathrm{d} k_1
        \mathrm{d} \x
        \mathrm{d} k_2 
        = \kappa(X,Y),
      \end{align*}
      where 
      \[
      G(k_1,\x,k_2) = f(k_1,\mathbf{x},k_2)\Psi(\x).
      \]

\subsubsection*{Positivity of $\kappa(X,Y)$} 
On $\mathcal{C}$, we have $\Psi(\x) \ge 0$ and thus $G(k_1,\x,k_2) \ge 0$, so it suffices to prove that 
      \[
       \int_{K} \int_{W} \int_{K}
          G(k_1,\x,k_2)
          \mathrm{d} k_1
        \mathrm{d} \x
        \mathrm{d} k_2 > 0
      \]
      for some $W \subset \mathcal{C}$. We show that we may take 
      \[
      W = \{
      \y \in H: | \mathbf{z} - \y | \le \mu
      \}, 
      \]
      where $\mu>0$ is some small positive parameter and
$       z_i = - v_i R + \delta_i.
 $
       The $\delta_i$ are chosen to be small perturbations such that $\sum_{i=1}^n \delta_i =0$ and 
      $z_i > z_{i+1}$ if $v_i = v_{i+1}$, and $R$ is a sufficiently large parameter. 
      We recall the definitions of $D_+(\x)$ and $D_-(\x)$ in \eqref{eq:defofD+} and \eqref{eq:defofD-}. 
      By \eqref{eq:basic_infty_norm_inequ}, for $\x \in W$ we have 
      \[
      | X k_1 D_+(\mathbf{x}) k_2| \le
      n^3 | X|  | D_+(\mathbf{x})| \leq
      n^3 | X| e^{-R + \delta_1 + \mu }
      \]
      and
      \[
      |  k_2^{-1} D_-(\mathbf{x}) k_1^{-1} Y|\le n^3 | Y| |   D_-(\mathbf{x})| \leq  n^3 | Y| e^{-R -\delta_n + \mu}.
      \]
      Thus we may choose $R$ sufficiently large and $\mu$ sufficiently small in such a way that $W$ lies inside the interior of $\mathcal{C}$ and $K \times W \times K$ is contained in the support of 
      \[
      \mathbf{1}_{[| X k_1 D_+(\mathbf{x}) k_2|\le 1]}
      \mathbf{1}_{[|  k_2^{-1} D_-(\mathbf{x}) k_1^{-1} Y|\le 1]}
      .
      \]
      Now let $ \nu >0$ be the minimum of the function $\Psi$ on the compact set $W$. Then clearly 
      \[
      c_n 
       \int_{K} \int_{W} \int_{K}
          G(k_1,\x,k_2)
          \mathrm{d} k_1
        \mathrm{d} \x
        \mathrm{d} k_2 \ge c_n  \, \nu \, \mu_H(  W )>0. \qedhere
      \]
\end{proof}

\end{document}
